\scanpage{28}
On \deleted{exige} \inserted{veut} de plus \deleted{compatibilité}, \deleted{\unsure{avec|que}} si $X/S$ est % L1
lisse et $Y\underset{i}{\subset}X$ est lisse sur $S$, de codim % L2
$d$ dans $X$, l'isom % L3
\[
\mathbb R i^!(1_X)\simeq\deleted{\mathbb Z}\mathbb Q_Y(-d) % L4
\]
compatible avec les isom déjà connus % L5
du point de vue $\ell$-adique .... % L6

6) \deleted{Chaque $\mathcal M(X)$ est plongé dans} % L7
Le foncteur % L8
\[
M\unsure{\leadsto|\mapsto}M(-1)=M\otimes\deleted{\mathbb Z}\mathbb Q(\unsure{+1|-1}) % L9
\]
\deleted{\unsure{est|de}} $\mathcal M^+(\unsure{X|A})$ dans lui-$=$ est \uline{pleinement} \uline{fidèle} % L10
mais pas une équivalence en général. % L11
Il y a donc une façon canonique % L12
d'élargir $\mathcal M^+(X)$ en \deleted{\unsure{une}} $\mathcal M(X)$ de telle façon % L13
que \deleted{\unsure{de}} $\otimes\deleted{\mathbb Z}\mathbb Q(-1)$ devienne une équivalence, % L14
\deleted{\unsure{ce}}, en prenant la \inserted{pseudo-}limite inductive % L15
% slanted
\margin{$\mathcal M^+(X)$ est une \uline{sous-catégorie abélienne épaisse de} $\mathcal M(X)$ ; [\unsure{l'appart.} à $\mathcal M^+(X)$ se vérifie fibre par fibre ----]} % L15a
des % L16
\diagram{
\begin{tikzcd}[ampersand replacement=\&]
\mathcal M^+ \arrow[r,"{\otimes\deleted{\mathbb Z}\mathbb Q(-1)}"] \& \mathcal M^+ \arrow[r,hook,"{\otimes\deleted{\mathbb Z}\mathbb Q(-1)}"] \& \mathcal M^+ % L16a
\end{tikzcd}
}
. On prolonge % L16b
$\mathcal M(X)$ la structure $\otimes$, alors $\deleted{\mathbb Z}\mathbb Q(-1)$ % L17
devient inversible, soit $\deleted{\mathbb Z}\mathbb Q(1)$ son inverse, et % L18
tt élément de $\mathcal M(X)$ peut s'écrire ss forme % L19

% archive sheet number: 27
